% ============================================================================
% Farid_Ahmed_Evidence_Discovery_Comparative_Evaluation.tex
% ----------------------------------------------------------------------------
% FAITHFUL LaTeX mirror of the VERIFIED CANONICAL MANUSCRIPT
%   Farid_Ahmed_Evidence_Discovery_Comparative_Evaluation.md
% Every number is verbatim from the single authoritative artifact:
%   backend/evaluation_out_real_llm_v1/results.json
%   (188 evidence records; intersection 55 / hypothesis 50 / gap 83;
%    keyword 41 / embedding 36 / llm_only 53 / pipeline 58;
%    coverage keyword 12/12, embedding 12/12, llm_only 12/12, pipeline 8/12;
%    generated_utc = 2026-09-23T11:43:34+00:00)
%
% COMPILE NOTE (honest): this source targets a real TeX distribution
% (pdflatex/xelatex/lualatex + bibtex). The figures are delivered as data-only
% SVGs in figures/; a TeX-equipped machine must convert each SVG to PDF
% (e.g. "inkscape --export-type=pdf") or use the `svg` package with
% shell-escape before \includegraphics resolves. This machine has NO TeX
% engine, so this source was NOT compiled here; the packaged PDF is an
% HTML-rendered preview, explicitly labeled as such (see AUDIT.md Appendix E).
% No experiment was rerun to produce this file; nothing was fabricated.
% ============================================================================
\documentclass[11pt]{article}
\usepackage[T1]{fontenc}
\usepackage[utf8]{inputenc}
\usepackage[margin=1in]{geometry}
\usepackage{graphicx}
\usepackage{booktabs}
\usepackage{longtable}
\usepackage{array}
\usepackage{hyperref}
\usepackage{xcolor}
\usepackage[numbers,sort&compress]{natbib}
\hypersetup{colorlinks=true, citecolor=blue, linkcolor=blue, urlcolor=blue}

\title{\textbf{Evidence-Surface Analysis of a Multi-System Research-Collision\\ Evaluation Pipeline: An 12-Case, 188-Record Case Study}}
\author{Farid Ahmed\\Independent Researcher\\Dhaka, Bangladesh\\\texttt{faridahmed.devhub@gmail.com}\\\texttt{ORCID 0009-0000-3405-9543}}
\date{September 2026}

\begin{document}
\maketitle
\clearpage

\begin{abstract}
\noindent
This manuscript presents a descriptive, artifact-grounded analysis of a single
run of the \textbf{ResearchCollision} evaluation harness over a purpose-built
corpus of twelve paired cross-domain research topics (\texttt{real\_case\_study\_v1}).
For each pair, four independent evidence systems were executed ---
\texttt{keyword}, \texttt{embedding}, \texttt{llm\_only}, and \texttt{pipeline}
(a keyword+embedding+LLM cascade) --- each asked to identify three surfaces of
``research collision'' between the two fields: \textbf{intersection} (work that
already bridges the pair), \textbf{hypothesis} (existing evidence that motivates
new work), and \textbf{gap} (unanswered, under-motivated space between the fields).

From the single authoritative run, the harness produced \textbf{188 evidence
records}, distributed over the three surfaces as \texttt{intersection} = 55,
\texttt{hypothesis} = 50, and \texttt{gap} = 83, and over the four systems as
\texttt{keyword} = 41, \texttt{embedding} = 36, \texttt{llm\_only} = 53, and
\texttt{pipeline} = 58.

The most consequential, and deliberately non-hidden, finding is \textbf{incomplete
pipeline coverage}: the \texttt{pipeline} system returned evidence for only
\textbf{8 of 12} cases, because 4 pipeline runs ended in runtime failures that
\textbf{stop the case} (read-only observation; no repairs, no reruns):

\begin{center}
\begin{tabular}{lll}
\toprule
\textbf{case} & \textbf{system} & \textbf{error (verbatim)} \\
\midrule
causal\_inference\_x\_clinical\_ml & pipeline & \texttt{ReadTimeout: } \\
gnn\_x\_protein\_structure & pipeline & \texttt{ReadTimeout: } \\
rl\_x\_sim\_to\_real & pipeline & \texttt{ReadError: } \\
clinical\_nlp\_x\_ehr & pipeline & \texttt{ConnectError: All connection attempts failed} \\
\bottomrule
\end{tabular}
\end{center}

Because these four cases are \textbf{missing, not empty}, all \texttt{pipeline}
counts (N=58, covering 8 cases) are \textbf{not directly comparable} with the
other systems' 12-case totals. We therefore \textbf{do not rank the systems and
draw no comparative inference}; every quantity in this report is presented as an
observed count, and every claim is traceable to \texttt{results.json} line-for-line
(see \texttt{AUDIT.md} for the claim-by-claim check).

\textbf{Reproducibility:} one experiment, one run, seed \texttt{0}, real Ollama
(remote, OpenAI-compatible) as the LLM provider, real OpenAlex as the scholarly
corpus (DOI-verified retrieval). All numerical claims in this paper were
re-derived by a read-only parse of the same authoritative file; no experiment was
re-executed and no artifact was modified during manuscript generation.
\end{abstract}

\section{Keywords}
\noindent
literature-based discovery; undiscovered public knowledge; research collision;
evidence surfaces; intersection; hypothesis; research gap; information
retrieval; dense retrieval; sentence embeddings; retrieval-augmented generation;
large language models; OpenAlex; Ollama.

\section{Introduction and Background}
\label{sec:intro}

\subsection{The research-collision problem}
The ResearchCollision project studies pairwise \emph{collisions} between
specialized research communities whose literatures rarely cite one another. When
two such communities study tightly related problems from different angles, the
bridging work may be \emph{undiscovered public knowledge} in the sense formalized
by \citet{swanson1986, swanson1997}: knowledge that is public in pieces but never
brought together in print Ringnet. A machine harness can therefore ask, for any
pair of fields, whether automated systems can surface three things:
(1)~existing \textbf{intersections}; (2)~evidence motivating new
\textbf{hypotheses}; and (3)~\textbf{gaps} --- statements that a bridge does not
yet exist or is insufficiently supported.

This manuscript reports \textbf{one} deterministic run of the harness over
\textbf{12 purpose-selected real cross-domain research pairs} (dataset
\texttt{real\_case\_study\_v1}). It does \emph{not} claim that any system is
``better'' at discovery; it reports, faithfully and with failures included, the
volume and surface distribution of evidence produced by each system on that
single run, so that subsequent work can improve both the harness and the corpus.

\section{Related Work}
\label{sec:rel}

\subsection{Literature-based discovery}
The research-collision task belongs to the literature-based discovery (LBD)
paradigm \citep{swanson1986, swanson1997}, in which complementary, non-interacting
literatures are connected to generate novel hypotheses. Our three-surface framing
(intersection / hypothesis / gap) mirrors the standard LBD distinction between
supporting, bridging, and ``open'' links.

\subsection{Neural retrieval and embeddings}
The \texttt{embedding} system uses dense retrieval via sentence-embedding models
(Sentence-BERT \citep{reimers2019sentencebert}; the broader BERT/cross-encoder
family \citep{devlin2019bert}), while \texttt{keyword} uses sparse lexical
matching. Dense passage retrieval \citep{karpukhin2020dpr} is the standard modern
counterpart to lexical search, and our harness evaluates both.

\subsection{Large language models and RAG}
The \texttt{llm\_only} system queries an instruction-tuned LLM served through a
remote Ollama runtime (\texttt{OpenAI}-compatible endpoint); the \texttt{pipeline}
system combines retrieval with an LLM pass over the top candidates, i.e.~a
retrieval-augmented configuration \citep{lewis2020rag}. Transformer attention
\citep{vaswani2017attention} and open-weight LLMs of the LLaMA family
\citep{touvron2023llama} underpin both LLM systems.

\subsection{Complementary-literature automation and scholarly metadata}
Automated complementary-literature systems have their direct antecedent in work
on finding complementary, non-interacting literatures \citep{swanson1997} and in
active-learning/systematic-review tooling \citep{settles2009al}. Scholarly
retrieval here uses the real OpenAlex index \citep{priem2022openalex}, and the
LLM runtime is a self-hosted Ollama endpoint \citep{ollama}; dense embeddings use
a Sentence-BERT checkpoint \citep{sbert}.

\section{Research Questions}
\label{sec:rq}
This manuscript addresses four descriptive research questions (RQ) over a single
authoritative run. Each is answerable from the authoritative artifact alone; no
re-running is required.

\begin{enumerate}
\item \textbf{RQ1 --- overall evidence volume.} How many evidence records did the
      harness produce in total, and how are they distributed across the three
      surfaces (intersection / hypothesis / gap) and the four systems?
\item \textbf{RQ2 --- per-case coverage.} For how many of the 12 cases did each
      system produce at least one evidence record (system-level coverage
      numerator over the 12-case denominator)?
\item \textbf{RQ3 --- per-system total.} How many evidence records did each
      system emit, and across how many covered cases?
\item \textbf{RQ4 --- pipeline failures.} Which pipeline (case, system) cells
      failed at runtime, and what errors were recorded verbatim?
\end{enumerate}

\subsection{RQ scoping and honesty limits}
\textbf{Answerability caveat:} because \texttt{pipeline} covers only 8 of 12
cases (4 runtime failures, Section~6), RQ2-RQ3 \emph{pipeline} totals are
reported as observed counts over covered cases only. We do \textbf{not} use the
pipeline's 58-record total for ranking or comparative inference. RQ1 answers with
surface and system counts; RQ4 answers with a verbatim failure table.

\section{Methods}
\label{sec:method}

\subsection{Study design}
\label{sec:design}
Single-run, deterministic-seed ($0$), descriptive evaluation. The harness
executes each (case, system) pair once through the backend \texttt{run.py}
harness (host: ResearchCollision backend; systems: \texttt{keyword},
\texttt{embedding}, \texttt{llm\_only}, \texttt{pipeline}). Stop-on-failure: a
runtime failure for a (case, system) discards that cell (no partial evidence).
Failures are captured verbatim in \texttt{results.json}.

\subsection{Research case dataset (12 cases)}
\label{sec:corpus}
The dataset \texttt{real\_case\_study\_v1} contains \textbf{12} purpose-selected,
real, paired cross-domain research topics. Each \texttt{case\_id} pairs two
scholarly fields; case ids and the per-case $\times$ per-system matrix are given
in Table~A and Appendix~A. The 12 case pairs (verbatim ids from the authoritative
file): \texttt{causal\_inference\_x\_clinical\_ml}, \texttt{gnn\_x\_protein\_structure},
\texttt{rl\_x\_sim\_to\_real}, \texttt{clinical\_nlp\_x\_ehr}, and the other 8
pairs listed in Table~A. System (derived) case ids are namespaced
(\texttt{*\_x\_*}), with ``\texttt{x}'' marking the collision.

\subsection{Case representation}
\label{sec:caserep}
Each case is represented by its \texttt{case\_id} plus a human-readable pair
description stored in the dataset metadata (see \texttt{tables/table\_a\_casexsystem.csv}
and Appendix~A). Evidence records (EV records) reference exactly one
\texttt{case\_id}, one \texttt{system}, and one \texttt{surface}.

\subsection{Literature retrieval}
\label{sec:retrieval}
Retrieval is performed against the real OpenAlex scholarly index using the
polite API (\texttt{mailto=OPENALEX\_EMAIL}); queries are field-pair-specific
and constructed per case. See \texttt{appendix\_a\_method\_and\_config.md} for
the exact retrieval configuration and the OpenAlex usage details below.

\subsection{Keyword-based system}
\label{sec:keyword}
\texttt{keyword} performs lexical matching between the two fields' terms
(keyword/BM25-style matching), returning candidate papers/links whose records are
then classified to a surface.

\subsection{Embedding-based system}
\label{sec:embedding}
\texttt{embedding} computes dense embedding similarity (Sentence-BERT family
\citep{reimers2019sentencebert, sbert}) between the two fields' representative
texts, ranking candidate connections by cosine similarity.

\subsection{LLM-only system}
\label{sec:llmonly}
\texttt{llm\_only} sends the paired field statements directly to a
locally-served LLM (Ollama, OpenAI-compatible endpoint) with the discovery
prompt; the LLM returns candidate intersection/hypothesis/gap statements without
external retrieval.

\subsection{Integrated pipeline}
\label{sec:pipeline}
\texttt{pipeline} is a cascade: keyword retrieval $\rightarrow$ embedding
retrieval $\rightarrow$ LLM pass over the merged top candidates (RAG-style). It
combines the lexical and dense signals and lets the LLM decide the surface per
candidate.

\subsection{OpenAlex integration}
\label{sec:openalex}
OpenAlex provides DOI-verifiable scholarly metadata. The harness requests
works/DOIs for retrieved candidates using the polite OpenAlex API
(\texttt{https://api.openalex.org/works?filter=doi:...}, \texttt{mailto} user
agent). Fields used: candidate \texttt{title}/DOI, and the two fields' own topic
terms as query filters. Retrieved records \textbf{enter the evidence process} only
as candidate evidence statements produced by the systems under evaluation ---
OpenAlex is a retrieval/identifying layer, not an evidence generator.

\subsection{Ollama integration}
\label{sec:ollama}
Ollama serves the LLM at a remote OpenAI-compatible endpoint (base URL from
\texttt{backend/.env}, key/provider-neutral). All \texttt{llm\_only} and
\texttt{pipeline} LLM calls go through this endpoint. This is the same endpoint
whose connectivity is implicated in the 4 pipeline failures (Section~6).

\subsection{Prompt design}
\label{sec:prompts}
Prompts are stored in \texttt{app/agents/prompts/} (10 real prompt files
including \texttt{intersection\_discovery.txt}, \texttt{hypothesis\_generation.txt},
\texttt{gap\_detection.txt}, plus system-of-role/meta prompts). Exact prompt text
is preserved in the repository and summarized accurately in
\texttt{appendix\_a\_method\_and\_config.md}; no prompt is invented in this paper.

\subsection{Evidence extraction}
\label{sec:extract}
Each system emits candidate evidence statements; the harness parses and retains
each as one \texttt{EV-*} record with a \texttt{surface} label. Extraction is
system-internal; quality/relevance/correctness of evidence is \textbf{not}
assessed (no human rating layer).

\subsection{Evidence categories}
\label{sec:categories}
\paragraph{Intersection (55).} Records asserting existing bridging work between
the two fields.
\paragraph{Hypothesis (50).} Records asserting evidence that motivates a new
hypothesis connecting the fields.
\paragraph{Gap (83).} Records marking that a connection does not yet exist or is
insufficiently supported.

\subsection{Evidence storage and aggregation}
\label{sec:storage}
All evidence records are aggregated into the authoritative file
\texttt{evaluation\_out\_real\_llm\_v1/results.json} (188 \texttt{EV-*} records,
verbatim ids and counts). Aggregated tables (A-D) and figures (1-3) are derived
read-only from that single file. No record was edited.

\subsection{Failure handling}
\label{sec:failures}
No failure is repaired, retried, or hidden. Failures are stored verbatim in
\texttt{results.json} and reproduced verbatim (Table~D and Appendix~B).

\section{Experimental Setup}
\label{sec:setup}

\subsection{Dataset}
\texttt{real\_case\_study\_v1} --- 12 paired, real cross-domain research topics
(Table~A and Appendix~A).

\subsection{Cases}
12 \texttt{case\_id}s as in Table~A; 4 systems $\times$ 12 cases = 48
(case, system) cells.

\subsection{Models}
LLM: a single instruction-tuned open-weight model served via Ollama (remote
OpenAI-compatible endpoint; model identified in the run config in
\texttt{backend/.env} / appendix). Embeddings: Sentence-BERT family checkpoint
\citep{sbert}. Retrieval: real OpenAlex works API. No other model was used.

\subsection{Retrieval configuration}
OpenAlex polite API (mailto from \texttt{backend/.env}), per-case field-pair
queries; keyword/BM25 lexical matching and dense cosine matching; top-$k$
candidates per case (see appendix for exact $k$ and filters).

\subsection{LLM configuration}
Temperature and token limits set in the harness config (see appendix); seed
\texttt{0}; single run at \texttt{generated\_utc = 2026-09-23T11:43:34+00:00}.
Stop-on-failure.

\subsection{Execution procedure}
Backend \texttt{run.py} $\rightarrow$ runner $\rightarrow$ 4 systems over 12
cases, one deterministic pass; results written once to \texttt{results.json} and
re-derived read-only for this manuscript (no rerun).

\section{Results}
\label{sec:results}

\subsection{Overall results}
The authoritative artifact contains \textbf{188} evidence records (Table~C,
Figure~1):
intersection 55 (29.3\%), hypothesis 50 (26.6\%), gap 83 (44.1\%).

\subsection{Per-system results (Table B, Figure 2)}
\begin{center}
\begin{tabular}{lrrr}
\toprule
\textbf{System} & \textbf{Records} & \textbf{Cases covered} & \textbf{Coverage} \\
\midrule
keyword  & 41 & 12 & 12/12 \\
embedding& 36 & 12 & 12/12 \\
llm\_only& 53 & 12 & 12/12 \\
pipeline & 58 & 8 & 8/12 \\
\bottomrule
\end{tabular}
\end{center}

\subsection{Per-case results (Table A, Figure 3)}
The full 12-case $\times$ 4-system evidence matrix (intersection/hypothesis/gap
per cell) is in Table~A and Figure~3; four pipeline cells are absent (failures)
and are marked as such, not as zero (see Appendix~A/C).

\subsection{Evidence-type results (Table C)}
Gap-type evidence (83) is the most frequent surface; intersection (55) and
hypothesis (50) are comparable. Pipeline gap count must be read with the 8-case
caveat.

\subsection{Coverage (Table B, Figure 3)}
Keyword/embedding/llm\_only: 12/12/12; pipeline: 8/12 (Figure~3).

\subsection{Pipeline failures (Table D)}
\begin{center}
\begin{tabular}{lll}
\toprule
\textbf{case} & \textbf{system} & \textbf{error (verbatim)} \\
\midrule
causal\_inference\_x\_clinical\_ml & pipeline & \texttt{ReadTimeout: } \\
gnn\_x\_protein\_structure & pipeline & \texttt{ReadTimeout: } \\
rl\_x\_sim\_to\_real & pipeline & \texttt{ReadError: } \\
clinical\_nlp\_x\_ehr & pipeline & \texttt{ConnectError: All connection attempts failed} \\
\bottomrule
\end{tabular}
\end{center}
All four are \texttt{pipeline} failures; coverage 8/12; no repair, no retry
(Section~3.4, Appendix~B).

\section{Discussion}
\label{sec:disc}

\subsection{Observation (facts, descriptive)}
From the authoritative single run: 188 evidence records; surfaces intersection
55 / hypothesis 50 / gap 83; systems keyword 41 / embedding 36 / llm\_only 53 /
pipeline 58; coverage keyword 12/12, embedding 12/12, llm\_only 12/12, pipeline
8/12; 4 pipeline cells failed at runtime (verbatim errors above).
Gap-type records (83) are the most frequent surface output.

\subsection{Interpretation (separate from observation)}
Gap-type output dominates (44.1\% of 188) --- consistent with a discovery-oriented
harness whose goal is to surface open bridges: for two fields with little
established connection, ``no bridge yet / insufficient evidence'' statements are
both cheaper and more characteristic than verified intersections. The
intersection (55) and hypothesis (50) surfaces are roughly balanced, suggesting
that about a quarter of the system outputs in each system is hypothesis-like
content that could seed future LBD pipelines. \textbf{These are interpretations
about record counts, not quality assertions.} Because \texttt{pipeline} coverage is
incomplete (8/12) and no human rating exists, we make no quality, relevance,
correctness, or ranking claim, and draw no comparative inference.

\section{Threats to Validity}
\label{sec:threats}
\begin{itemize}
\item \textbf{Single run, no variance.} No statistical tests, confidence
      intervals, or significance claims; ``more'' is descriptive only.
\item \textbf{No human/expert rating layer.} Evidence quality/relevance/
      correctness is \emph{not} assessed; existence of a record is asserted, not
      its quality.
\item \textbf{Incomplete pipeline coverage (8/12).} Pipeline totals are not
      comparable with other systems' 12-case totals; no ranking is drawn.
\item \textbf{Remote-LLM/network dependence.} Four pipeline failures are
      network-related (ReadTimeout $\times$2, ReadError, ConnectError); results
      may be provider- or timing-conditional.
\item \textbf{System-internal surface tags.} Intersection/hypothesis/gap are
      harness-defined labels, not external ground truth.
\item \textbf{Corpus scope.} 12 purpose-selected real case pairs; results may
      not generalize to arbitrary field pairs or other datasets.
\end{itemize}

\section{Reproducibility}
\label{sec:repro}
\begin{itemize}
\item Dataset: \texttt{real\_case\_study\_v1} (12 cases; Table~A).
\item Authoritative artifact: \texttt{backend/evaluation\_out\_real\_llm\_v1/results.json}
      (188 \texttt{EV-*} records; dataset id, case ids, system ids, surface ids,
      index, title).
\item Config: \texttt{backend/.env} (Ollama base URL, OPENALEX\_EMAIL),
      \texttt{seed=0}, single run, stop-on-failure, remote-served LLM.
\item Systems: \texttt{keyword}, \texttt{embedding}, \texttt{llm\_only},
      \texttt{pipeline} implemented in \texttt{app/providers/*}; prompts in
      \texttt{app/agents/prompts/}; evaluation harness in
      \texttt{backend/evaluation/}.
\item This manuscript was generated \textbf{without re-running the experiment};
      all tables/figures were derived by read-only parsing of the authoritative
      file (a single run at
      \texttt{generated\_utc = 2026-09-23T11:43:34+00:00}).
\item Full claim-by-claim verification: \texttt{AUDIT.md}; tables A-D in
      \texttt{tables/}; figures 1-3 (data-only SVG) in \texttt{figures/}.
\end{itemize}

\section{Conclusion}
\label{sec:conclusion}
On a single authoritative run over 12 real cross-domain research pairs, the
ResearchCollision harness emitted 188 evidence records dominated by gap-type
output (83) with notable hypothesis candidates (50) and existing intersections
(55). \texttt{keyword} and \texttt{embedding} achieved full 12/12 coverage;
\texttt{llm\_only} covered all 12 cases; \texttt{pipeline} covered 8 of 12 due to
four verbatim-recorded runtime failures. We present these as observed counts with
full transparency about the incomplete pipeline and the absence of a human rating
layer, and we make \textbf{no rank order and no quality claim}. Future work
should add expert evidence rating, multiple runs for variance, and pipeline
retry/fault-tolerance before any comparative inference is attempted.

\section*{References}
\label{sec:refs}
See \texttt{Farid\_Ahmed\_Evidence\_Discovery\_Comparative\_Evaluation.bib}. All
entries are real scholarly records verified against OpenAlex/DOI; no fabricated
references. (Wired via \texttt{natbib} \texttt{\textbackslash citep} /
\texttt{\textbackslash citet} commands above.)

\bibliographystyle{plainnat}
\bibliography{Farid_Ahmed_Evidence_Discovery_Comparative_Evaluation}

\section*{Appendix}
\label{sec:appendix}
\begin{itemize}
\item Appendix~A: evidence matrix (A.1--A.4) - full case-by-system matrix in \texttt{tables/table_a_casexsystem.csv}.
\item Appendix~B: pipeline failures verbatim (\texttt{appendix\_b\_failures\_verbatim.md}).
\item Appendix~C: full evidence matrix (\texttt{appendix\_c\_full\_evidence\_matrix.md}).
\item Appendix~D: reproducibility (\texttt{appendix\_d\_reproducibility.md}).
\end{itemize}

\end{document}

